Pour identifier un petit objet métallique, une élève mesure sa masse : \qty{\fpeval{round(\v{rho}*\v{dV},1)}}{g}. Elle verse ensuite \qty{\v{Vi}}{mL} d'eau dans un cylindre gradué et y plonge l'objet : le niveau monte à \qty{\fpeval{\v{Vi}+\v{dV}}}{mL}.

\begin{center}
\begin{tabular}{lc}
    \hline
    Métal & Masse volumique \\
    \hline
    Aluminium & \qty{2.70}{g/mL} \\
    Fer & \qty{7.87}{g/mL} \\
    Cuivre & \qty{8.96}{g/mL} \\
    Plomb & \qty{11.34}{g/mL} \\
    \hline
\end{tabular}
\end{center}

\begin{enumerate}[label=\alph*)]
    \item Quel est le volume de l'objet ? Comment s'appelle cette méthode de mesure ?
    \item Calculez la masse volumique de l'objet.
    \item À l'aide du tableau, identifiez le métal. Pourquoi la masse volumique permet-elle d'identifier une substance, alors que la masse ne le permet pas ?
\end{enumerate}
